%!TEX root = ../../main.tex
%% ===============================================================
%% 6.3 사건 특징 ($F_{\mathrm{ev
%% 파일: chapters/06_features/03_event_features.tex
%% 상위: chapters/06_features/chapter.tex
%% 정의 라벨: sec:feat-event
%% 참조 라벨: -
%% 포함 객체: -
%% 인용: -
%% 상태: 초안 (docs/OUTLINE.md 참고)
%% ===============================================================

\section{사건 특징 ($F_{\mathrm{ev}} = 44$)와 사건 토큰}
\label{sec:feat-event}

구간 $[b_0, b_1)$ 안에서 발생한 사건을 종류별·팀별로 계수하여 $x^{\mathrm{ev}}_\ell \in \R^{44}$를 만든다. 22 종의 사건(킬, 현상금, 셧다운, 연속 킬, 다중 킬, 에이스, 용, 바론, 전령, 아타칸, 유충, 포탑, 억제기, 방패, 오브젝트 현상금, 와드 설치·제거, 제어 와드 설치·제거, 아이템 구매·판매·취소)에 두 팀을 곱한 것이다. 교차 주의 모델을 위해서는 집계 대신 최대 $K = 64$개의 이산 토큰 $E \in \R^{K \times D_e}$를 별도로 구성한다. 각 토큰은 종류 해시, 행위자 ID, 팀과 함께 상대 시각, 온셋까지의 시간, 위치, 값, 중요도 사전값 등 $D_e = 12$개의 연속 특징을 갖는다.
